% ===========================================================================
\section{Discussion}
\label{sec:discussion}
% ===========================================================================

\subsection{What the architecture actually bought}

The claim in Section~\ref{sec:architecture} is testable against the project's own history,
and it is worth being honest about which parts of it paid off.

The clearest return was on \emph{late additions}. Fugitives and cement were built last, after
the analysis layer was mature. Neither required any change to the uncertainty analysis, the
key-category screen, the benchmark comparison, the provincial allocation or the figure
generation; both appeared in every national output on first registration. In a spreadsheet
inventory, adding a sector means extending every downstream sheet by hand, and the cost of a
new sector rises with the number of sectors already present. Here it is flat.

The second return was on \emph{invariants catching real errors}. Several of the \num{27}
contract tests exist because they failed, and Table~\ref{tab:qa} is the audit rather than the
anecdote: each row is a defect this project shipped, what it would have cost, how it was
found, and the guard that now makes it a build failure. The pattern across them is worth more
than any single row. \textbf{Every one was invisible to review}, because nothing about the
resulting numbers looked wrong---and the later ones were invisible to the numeric tests too,
because they were claims about words, boundaries or layout rather than about digits.

\begin{table*}[htbp]
\centering
\footnotesize
\setlength{\tabcolsep}{4pt}
\begin{tabular}{@{}p{0.19\linewidth}p{0.17\linewidth}p{0.16\linewidth}p{0.18\linewidth}p{0.20\linewidth}@{}}
\toprule
\textbf{Defect} & \textbf{Consequence} & \textbf{Detection} & \textbf{Correction} &
\textbf{Regression guard} \\
\midrule
Two sectors proposing one category name & silent double-count & contract test on
registration & globally unique names & category-name invariant \\
Correlation group spanning two categories & variance apportionment ill-defined & contract
test & groups within one category & correlation-group invariant \\
Wrong parallel factor table & category trace gases 8\% out, \COtwo{} exact & factor audit &
category-specific factor identity & factor-identity test \\
Ambiguous integer year labels & mixed-period national total & period audit & explicit
start/end dates & period-schema test \\
Unread public activity table & \SI{3.40}{\mega\tonne} coverage gap & repository audit & OCAC
liquid customer class read & completeness test \\
Retired provincial power proxy & \SI{8.8}{\mega\tonne} in the wrong province & spatial
diagnostic & plant-wise generation & allocation test \\
Source documents never archived & provenance unverifiable & archive audit & source-document
registry & document-manifest test \\
Figure sources named by year & two manuscripts shipped pre-harmonisation figures & external
review & year resolved at read time, absence raises & figure-source test \\
Stale gas decomposition & a decomposition summing to \SI{455.14}{\mega\tonne} against a
printed \SI{448.00}{} & external review & generated category~$\times$~gas ledger &
gas-total invariant \\
Declared-debt list & printed quantities tracked but unverified & external review & the list
deleted; no exemption without a reason & unconditional coverage test \\
Two category hierarchies & \num{20} categories screened against \num{24} reported & external
review & IPCC category made canonical & one-hierarchy test \\
Stale comparative language & ``doubles'' surviving a ratio of \num{1.61} & external review &
claims registry binding phrase to relation & semantic-claims test \\
One correlation group per row & a factor correlation reported as widening a category by more
than half when it moves it by nothing & external review & uncertainty carried per parameter &
decomposition-identity test \\
Content printed off the page & \SI{104}{pt} of a table outside the margin & external review &
tables narrowed & log-derived layout test \\
\addlinespace
Superseded output vintage & a retired FY2019--20 fugitives file quoting
\SI{3.81}{\mega\tonne} beside the then-current \SI{24.69}{} (the envelope's top) & rebuild
audit & vintage retired &
orphan and twin sweep over \texttt{outputs/} \\
Staleness guard reading it & the shipped-versus-engine check ran on three sectors of nine and
reported passing & rebuild audit & the check reads each sector's declared categories &
per-sector coverage asserted, not assumed \\
Orchestrator declaring outputs it did not write & four phases named files that existed
nowhere and printed OK beside them & rebuild audit & the declared artefact must be written
during the run & orchestrator output check \\
Run manifest listing a directory scrape & a sector's manifest claimed any leftover file as
its own product & rebuild audit & the manifest records what the run wrote & write-tracked
manifest \\
The build defined three times & the branch that gates merges built a smaller inventory than
either local command & rebuild audit & the three lists reconciled & build-definition test \\
The release record & it printed \SI{455.14}{\mega\tonne} and described a mixed-year
inventory three rounds after both stopped being true & rebuild audit & repaired and anchored
& the release record is a claims document \\
\bottomrule
\end{tabular}
\caption{Defects this inventory shipped, and what now prevents each. The right-hand column is
the deliverable: a defect without a regression guard has been corrected, not closed. The
detection column is the uncomfortable one---the middle entries were found by an external
reader rather than by the suite, and each of those had passed a suite of more than a thousand
tests, because the property that was wrong had never been made checkable. The last six were
found by rebuilding the whole release from one commit and asking a narrower question than the
suite asks: not whether each artefact is correct, but whether anything still writes it. Four
of the six are the same defect wearing different clothes---a check that reports passing while
testing nothing---which is the failure mode a growing test suite is least able to see in
itself.}
\label{tab:qa}
\end{table*}

The third, and least anticipated, was that \emph{closures turned prose caveats into
arithmetic}. The four enforced closures in Section~\ref{sec:closures} existed first as written
warnings---``this biases the sector high'', ``there is no national figure to remove them
by''---in the documentation of sectors that had done the honest thing and declared a known
defect. Building the counterpart sector converted each warning into a number. That is a
general pattern worth naming: a well-documented known defect is a specification for the work
that fixes it.

What the architecture did \emph{not} buy is accuracy. The engine's activity data are still
national statistics of variable quality, its factors are still mostly IPCC Tier~1 defaults,
and the largest sector rests on a livestock census. Structure prevents inconsistency; it does
not manufacture information.

\subsection{Two undeclared assumptions worth a table apiece}

Two results in this paper share a shape, and the shape is the point.

Sharing parameters by physical identity across modules, rather than confining them within
one, leaves the national uncertainty at $\pm\SI{11.72}{\percent}$ and the
road-transport category unmoved---a result that reverses
what the previous edition reported, for reasons Section~\ref{sec:uncertainty} sets out. The sewer methane correction factor
moves domestic wastewater from \SI{15.39}{\mega\tonne} to \SI{24.46}{\mega\tonne}---a factor of \num{1.6} on a category that is \SI{5.5}{\percent} of the national inventory. In both
cases the assumption is defensible in either direction, the evidence does not settle it, and
the conventional treatment is to pick one and not mention it.

We publish both endpoints in both cases. That is not fence-sitting; it is the recognition
that a single number carries a false claim about the state of knowledge, and that the correct
output of an analysis whose answer depends on an unsettled assumption is the dependence
itself. In both cases we also state what would settle it: for correlation, a Pakistani fuel
carbon-content measurement programme; for the MCF, dissolved-methane sampling on urban drains
in two or three cities.

\subsection{Implications for Pakistani inventory work}

Three practical conclusions follow, in descending order of confidence.

\textbf{Resolve gas-system practices before refining the vehicle fleet.} The 2019 segment
ledger now includes utility throughput and pipeline length, making fugitives
\SI{5.3}{\percent} of national emissions and \SI{3.3}{\percent} of variance. Road transport
is \SI{9.0}{\percent} and \SI{0.9}{\percent}. The binding fugitive evidence is therefore
pipeline material and pressure, LDAR, compressor configuration, processing throughput and
well status rather than another aggregate production factor.

\textbf{Industrial fuel combustion is the binding constraint on completeness.} At roughly
\SI{16}{\mega\tonne} still missing, the unestimated remainder alone would rank sixth among
this inventory's modules---just below residential fuel---while the complete \SI{43.98}{\mega\tonne}
category EDGAR reports would rank third, above road transport. Its absence is
most of why this inventory reports \SI{87}{\percent} of EDGAR's aggregate magnitude. No amount of
refinement elsewhere compensates for the unmeasured fuels within this partially estimated
industrial category.

\textbf{Methane is the metric-sensitive majority.} At \SI{54.6}{\percent} of the national
total on AR5, Pakistan's inventory is dominated by a short-lived gas, and mitigation
assessments will be unusually sensitive to the metric and time horizon chosen. Reporting both
AR5 and AR6 at every level of aggregation, as here, is a minimum; a country in this position
has a stronger case than most for also reporting on a short-horizon metric.

\subsection{Co-emitted air pollutants: a companion product, not a result of this paper}

Seven of the ten modules also compute an air-pollutant layer from the same activity
quantities, through EMEP/EEA factors \citep{emep2019,emep2023soils} rather than from a
separate study. That shared activity basis is the methodologically interesting part, and it
supports one robust statement: on primary \PM{}, household solid fuel is the largest estimated
source in this inventory by a wide margin, while being a minor greenhouse-gas category---the
two inventories rank sources close to inversely.

\textbf{It is reported in the supplement rather than here, deliberately.} The seven layers are
not of comparable assurance: two are bounds or screening estimates, transport is exhaust-only,
and industrial combustion and IPPU have no layer at all. A ranking drawn across evidence that
uneven is easy to over-read, and it should not carry weight the greenhouse-gas product has
earned and it has not. Supplementary Note~S1 gives the layer-by-layer assurance matrix, the
magnitudes and the figure, released as \textbf{PakEmis-Air beta}.

\subsection{Generalisability}

The architecture is not Pakistan-specific. The contract in Section~\ref{sec:architecture}
encodes properties of IPCC inventories generally---categories, gases, notation keys,
allocation proxies, correlation groups---and the accounting closures are instances of a
pattern (one national quantity consumed by two sectors) that recurs in every national
inventory: refinery fuel between 1A1b and 1A2, waste between disposal and incineration, land
between cropland and grassland.

The parallel-table trap of Section~\ref{sec:factors} is entirely general and, we would
expect, common. Any inventory compiled by a small team under time pressure, using the 2006
Guidelines, is exposed to it, and the exposure is invisible to \COtwo{}-based validation.
Checking which of the four Chapter~2 tables and which of the Chapter~3 tables each fuel-using
category actually drew from would be a cheap and, we suspect, productive audit of many
published national inventories.
